% P15.27: the remainder of Durrett Section 3.3 and its exercise presentations.
\chapter{Characteristic-Function Criteria and the Moment Problem}
\label{sec:p15-characteristic-function-criteria-and-moments}

This chapter completes the remaining constructions and moment arguments
of Durrett's Section~3.3, including its outstanding exercises
\cite[pp.~131--143]{DurrettPTEAuthor2019}.
The preceding chapter supplies characteristic functions, inversion,
the continuity theorem and moment derivatives. Here every probability
law is a Borel probability on \(\mathbb R\), every random variable is
finite and real unless another domain is explicitly stated, and
\(\varphi_X(t)=\mathbb E e^{itX}\).
The normal density and CDF are denoted by \(\phi,\Phi\).
All logarithms are natural.

\section{Approximation and integration premises}
\label{sec:p15-criteria-approximation}

\begin{hbtheorem}{thm:p15-criteria-real-integration}{proved=annex:p15.criteria-real-integration}{Interval approximation, substitution and Riemann--Lebesgue}
Every Borel set of finite Lebesgue measure can be approximated,
in measure of symmetric difference, by finite unions of bounded
intervals. Every complex \(g\in L^1(\mathbb R)\) can be approximated
in \(L^1\) by finite interval step functions and by continuous
functions of compact support. Consequently
\[
 \int_{\mathbb R}g(x)e^{itx}\,dx\longrightarrow0
                                         \quad(|t|\to\infty).
\]
In particular a law with a Lebesgue density has a characteristic
function tending to zero at infinity.

If \(T:I\to J\) is an increasing \(C^1\) bijection between open
intervals with \(T'>0\), or a decreasing one with \(T'<0\), then
for every nonnegative Borel \(h\),
\[
 \int_J h(y)\,dy=\int_I h(T(x))|T'(x)|\,dx.
\]
The same formula holds for complex integrable \(h\), with absolute
integrability equivalent on the two sides.
\end{hbtheorem}

\begin{hbtheorem}{thm:p15-criteria-weak-integration}{proved=annex:p15.criteria-weak-integration}{The source weak-integration and finite-approximation premises}
If \(\mu_n\Rightarrow\mu\) and \(g\geq0\) is continuous, then
\(\int g\,d\mu\leq\liminf_n\int g\,d\mu_n\).
If also \(\mu_n\Rightarrow\mu\), \(g,h\) are real continuous
functions, \(g>0\),
\(|h(x)|/g(x)\to0\) as \(|x|\to\infty\), and
\(\sup_n\int g\,d\mu_n\leq C<\infty\), then \(h\) is
integrable for all these laws and
\(\int h\,d\mu_n\to\int h\,d\mu\).

For every CDF \(F\), let \(\mu\) be its probability law.
For integers \(m\geq1\), let
\[
 x_{m,j}=Q_F((j-\tfrac12)/m),\quad
 \mu_m=\frac1m\sum_{j=1}^m\delta_{x_{m,j}}
                           \qquad(1\leq j\leq m).
\]
Its CDF satisfies \(\sup_x|F_m(x)-F(x)|\leq1/(2m)\);
thus \(\mu_m\Rightarrow\mu\). Repeated atoms are allowed.
If \(\mu((0,\infty))=1\), all these atoms are strictly positive.
These statements supply the full needed content of source
Exercises~3.2.4--3.2.5 and 3.2.10.
\end{hbtheorem}

\begin{hbdefinition}{def:p15-criteria-complex-log}{A fixed complex logarithm near one}
For \(|z|<1\), define
\[
 L(1+z)=\sum_{k=1}^\infty\frac{(-1)^{k+1}z^k}{k}.
\]
It is the continuous logarithm on this disk with \(L(1)=0\).
The powers and exponential use the preceding chapter's complex
conventions.
\end{hbdefinition}

\begin{hblemma}{thm:p15-criteria-complex-log}{proved=annex:p15.criteria-complex-log}{The logarithm and complex-power limits}
The series defining \(L\) converges absolutely and locally
uniformly, \(e^{L(w)}=w\), and for \(|z|\leq1/2\),
\[
 |L(1+z)-z|\leq |z|^2,\qquad |L(1+z)|\leq2|z|.
\]
If \(c_n\to c\in\mathbb C\), then
\((1+c_n/n)^n\to e^c\).
If continuous functions \(A,B\) on an interval have
\(e^A=e^B\) and agree at one point, then \(A=B\)
provided their difference takes values in \(2\pi i\mathbb Z\);
in particular continuous logarithmic lifts with the same
initial value agree.
\end{hblemma}

\section{The remaining convergence and series exercises}
\label{sec:p15-criteria-convergence-exercises}

\begin{hbtheorem}{thm:p15-criteria-tight-transforms}{proved=annex:p15.criteria-tight-transforms}{Equicontinuity and compact uniform convergence}
For a family of probability laws \((\mu_i)_{i\in I}\) satisfying
\(\sup_{i\in I}\mu_i(\{|x|>M\})\to0\) as \(M\to\infty\),
called a tight family,
their characteristic functions are uniformly equicontinuous:
for every \(\varepsilon>0\) there is \(\delta>0\) such that
\[
 |\varphi_i(t+h)-\varphi_i(t)|<\varepsilon
                 \quad(i\in I,\ t\in\mathbb R,\ |h|<\delta).
\]
If \(\mu_n\Rightarrow\mu\), the transforms converge uniformly
on every compact real interval. They need not converge uniformly
on all of \(\mathbb R\): for \(\delta_{1/n}\Rightarrow\delta_0\),
\(\sup_t|e^{it/n}-1|=2\).
This is source Exercise~3.3.13.
\end{hbtheorem}

\begin{hbtheorem}{thm:p15-criteria-uniform-sums}{proved=annex:p15.criteria-uniform-sums}{The full uniform-sum inversion and polynomial density}
If \(U_1,\ldots,U_n\) are independent uniform variables on
\((-1,1)\), with integer \(n\geq2\), their sum has the
continuous density
\[
 f_n(x)=\frac1\pi\int_0^\infty
                  \left(\frac{\sin t}{t}\right)^n\cos(tx)\,dt
 =\frac1{2^n(n-1)!}\sum_{j=0}^n(-1)^j\binom nj
                              (x+n-2j)_+^{n-1}.
\]
The integral is absolutely convergent, with \(\sin t/t=1\)
at zero. The density vanishes outside \([-n,n]\) and is
a polynomial of degree at most \(n-1\) on every interval
\((k,k+1)\), \(k\in\mathbb Z\).
This completes source Exercise~3.3.5.
\end{hbtheorem}

\begin{hbtheorem}{thm:p15-criteria-limit-operations}{proved=annex:p15.criteria-limit-operations}{Normal parameters, independent sums and infinite products}
The following source Exercises~3.3.7--3.3.9 hold.
\begin{enumerate}
\item If centered normal laws of variances \(v_n\geq0\)
converge weakly to a proper law, then \(v_n\to v<\infty\)
and the limit is the centered normal law of variance \(v\).
\item If \(X_n,Y_n\) are independent within each row, including
the limiting row, and \(X_n\Rightarrow X_\infty\),
\(Y_n\Rightarrow Y_\infty\), then
\(X_n+Y_n\Rightarrow X_\infty+Y_\infty\).
No independence between different rows is required.
\item If independent \(X_j\) have partial sums
\(S_n\to S_\infty\) almost surely, then
\[
 \varphi_{S_\infty}(t)=\lim_{n\to\infty}\prod_{j=1}^n\varphi_{X_j}(t)
                                  \qquad(t\in\mathbb R).
\]
This limit defines the infinite product; zero limits are allowed.
\end{enumerate}
\end{hbtheorem}

\begin{hbtheorem}{thm:p15-criteria-series-examples}{proved=annex:p15.criteria-series-examples}{The binary product and the Cantor transform}
For independent fair signs \(\varepsilon_j\),
\(\sum_{j\geq1}\varepsilon_j2^{-j}\) has the uniform law on
\((-1,1)\), and hence
\[
 \frac{\sin t}{t}=\prod_{j=1}^\infty\cos(t/2^j)
                                      \quad(t\in\mathbb R),
\]
with the value one at zero.
For independent fair \(0,1\) variables \(B_j\), set
\[
 C=2\sum_{j=1}^\infty B_j3^{-j}.
\]
Its law has no atoms and is supported on a Borel set of
Lebesgue measure zero. Its transform is
\[
 \varphi_C(t)=e^{it/2}\prod_{j=1}^\infty\cos(t/3^j),
 \qquad
 \varphi_C(3^k\pi)=i\prod_{j=1}^\infty\cos(\pi/3^j)\ne0
                                     \quad(k\geq0).
\]
Thus absence of atoms does not imply decay of a characteristic
function at infinity. These are source Exercises~3.3.10--3.3.11.
\end{hbtheorem}

\begin{hbtheorem}{thm:p15-criteria-moment-identities}{proved=annex:p15.criteria-moment-identities}{Gaussian moments and the absolute-moment integral}
A standard normal \(Z\) has all absolute moments finite, odd
moments zero, and
\[
 \mathbb EZ^{2n}=\frac{(2n)!}{2^n n!}=(2n-1)!!
                                  \quad(n\geq0),
\]
where \((-1)!!=1\).
For every real probability law, with infinite values permitted,
\[
 \mathbb E|X|=\frac2\pi\int_0^\infty
                    \frac{1-\operatorname{Re}\varphi_X(t)}{t^2}\,dt.
\]
These are source Exercises~3.3.12 and 3.3.15.
\end{hbtheorem}

\begin{hbtheorem}{thm:p15-criteria-local-zero}{proved=annex:p15.criteria-local-zero}{A characteristic-function neighborhood detects the zero limit}
For laws of \(Y_n\), one has \(Y_n\Rightarrow0\) if and only if
there is a real \(\delta>0\) such that
\(\varphi_{Y_n}(t)\to1\) for every \(|t|\leq\delta\).
In that case \(\mathbb P(|Y_n|>\varepsilon)\to0\) for
every \(\varepsilon>0\).
This is source Exercise~3.3.17.
\end{hbtheorem}

\begin{hbtheorem}{thm:p15-criteria-independent-series}{proved=annex:p15.criteria-independent-series}{Independent partial sums: weak, probability and almost-sure limits}
For independent real \(X_j\), if
\(S_n=\sum_{j=1}^nX_j\) converges in distribution to a
proper law, then \(S_n\) converges in probability and
almost surely to a finite variable with that law.
The proof supplies probability-Cauchy completeness and
the required first-crossing inequality: with
\(S_{m,j}=\sum_{r=m+1}^jX_r\), integers \(0\leq m<n\),
and \(a>0\),
\[
 \mathbb P\!\left(\max_{m<j\leq n}|S_{m,j}|>2a\right)
 \min_{m<k\leq n}\mathbb P(|S_{k,n}|\leq a)
           \leq\mathbb P(|S_{m,n}|>a).
\]
An empty sum is zero. This completes Exercise~3.3.18
and the actual needed content of Exercises~2.5.9--2.5.10.
\end{hbtheorem}

\begin{hbtheorem}{thm:p15-criteria-pitman-law}{proved=annex:p15.criteria-pitman-law}{The full derivative criterion for the constant weak law}
Let \(X_j\) be i.i.d. with characteristic function \(\varphi\),
and let \(S_n=\sum_{j=1}^nX_j\).
For a real \(a\),
\[
 S_n/n\longrightarrow a\ \text{in probability}
 \quad\Longleftrightarrow\quad
 \varphi'(0)=ia.
\]
Under either condition,
\(\varphi(t/n)^n\to e^{iat}\) for every real \(t\).
The derivative is with respect to a real parameter and is
two-sided. Its necessity proof resolves the possible
integer multiples of \(2\pi i\) in taking logarithms.
It need not imply \(\mathbb E|X|<\infty\).
For example, the normalized law
\[
 \mathbb P(X=(-1)^k k)=\frac{c}{k^2\log k},\quad k\geq2,
 \qquad c^{-1}=\sum_{k=2}^\infty\frac1{k^2\log k},
\]
has infinite first absolute moment and a finite derivative
\(\varphi'(0)=ia\), where
\(a=c\sum_{k=2}^\infty(-1)^k/(k\log k)\).
This supplies the full corrected Exercise~3.3.14 and
its referred example.
\end{hbtheorem}

\section{Polya's criterion and stable constructions}
\label{sec:p15-criteria-polya}

\begin{hblemma}{thm:p15-criteria-convex-derivative}{proved=annex:p15.criteria-convex-derivative}{The full convex-derivative premise}
Let \(f:(0,\infty)\to[0,\infty)\) be finite, convex and
nonincreasing, with \(f(t)\to0\) at infinity.
Its finite right derivative \(d(t)=f'_+(t)\) exists for every
\(t>0\), is nonpositive, nondecreasing and right-continuous,
and tends to zero at infinity. For \(0<a<b<\infty\),
\[
 f(b)-f(a)=\int_a^b d(s)\,ds.
\]
There is a sigma-finite Borel measure \(\eta\) on \((0,\infty)\)
such that \(\eta((a,b])=d(b)-d(a)\).
The derivative at zero is not required to be finite.
\end{hblemma}

\begin{hbtheorem}{thm:p15-criteria-polya}{proved=annex:p15.criteria-polya}{Durrett 3.3.22: representation and the source approximation route}
Suppose \(\varphi:\mathbb R\to[0,\infty)\) is even,
\(\varphi(0)=1\), nonincreasing and convex on \((0,\infty)\),
and
\(\lim_{t\downarrow0}\varphi(t)=1\),
\(\lim_{t\to\infty}\varphi(t)=0\).
Then a probability measure \(\nu\) on \((0,\infty)\) satisfies
\[
 \varphi(t)=\int_0^\infty(1-|t|/s)^+\,\nu(ds).
\]
In particular \(\varphi\) is a characteristic function.
If \(\varphi\) has finitely many linear pieces on the positive
half-line, \(\nu\) is supported on its positive slope-change
points \(s_j\), with masses \(s_j(d(s_j)-d(s_j-))\),
where \(d=\varphi'_+\).
The proof constructs \(\nu(ds)=s\,\eta(ds)\), proves its
normalization and uses the source's finite-support
approximation followed by the continuity theorem.
A separate direct product-space construction gives the same
existence conclusion.
If instead the limit at infinity is \(c\in[0,1)\), the
characteristic-function conclusion still holds, with representation
\[
 \varphi(t)=c+(1-c)\int_0^\infty(1-|t|/s)^+\,\nu_c(ds)
\]
for a probability \(\nu_c\) on \((0,\infty)\).
The resulting law has mass exactly \(c\) at zero.
For \(c=1\), \(\varphi\equiv1\) and the law is \(\delta_0\).
Continuity at zero cannot be omitted: \(\mathbf1_{\{0\}}\)
satisfies the other zero-tail conditions but is not a
characteristic function.
\index{Polya's criterion}
\end{hbtheorem}

\begin{hblemma}{thm:p15-criteria-binomial-series}{proved=annex:p15.criteria-binomial-series}{The full binomial-series premise}
For every real \(\beta\) and real \(z\) with \(|z|<1\),
\[
 (1-z)^\beta=\sum_{n=0}^\infty(-1)^n\binom\beta n z^n,
 \qquad
 \binom\beta n=\frac{\beta(\beta-1)\cdots(\beta-n+1)}{n!}.
\]
The term at \(n=0\) is one.
For \(0<\beta<1\), the numbers
\(c_n=(-1)^{n+1}\binom\beta n\), \(n\geq1\),
are positive and sum to one; the identity
\(1-(1-z)^\beta=\sum_{n\geq1}c_nz^n\) extends to
\(-1\leq z\leq1\) with absolute convergence.
\end{hblemma}

\begin{hbtheorem}{thm:p15-criteria-stable-functions}{proved=annex:p15.criteria-stable-functions}{Both source constructions of exp(-abs(t) to a power)}
For \(0<\alpha\leq2\), \(e^{-|t|^\alpha}\) is a characteristic
function. For \(0<\alpha\leq1\), Polya's criterion supplies
the first source proof. For \(0<\alpha<2\), the separate
Spitzer construction uses
\[
 \psi(t)=1-(1-\cos t)^{\alpha/2},\qquad
 e^{-|t|^\alpha}
 =\lim_{n\to\infty}
       \{\psi(\sqrt2\,t\,n^{-1/\alpha})\}^n.
\]
The endpoint \(\alpha=2\) is the centered normal law of variance
two; \(\alpha=1\) is the standard Cauchy law.
For \(\alpha>2\), this function is not a characteristic function.
For independent variables with this transform,
\[
 (X_1+\cdots+X_n)/n^{1/\alpha}\ \stackrel d=\ X_1.
\]
This includes Example~3.3.23 and Exercises~3.3.19--3.3.20.
\end{hbtheorem}

\section{Agreement sets and cancellation}
\label{sec:p15-criteria-agreement}

\begin{hbtheorem}{thm:p15-criteria-periodic-example}{proved=annex:p15.criteria-periodic-example}{Durrett 3.3.24: the complete periodic Fourier calculation}
Let \(\psi\) be the continuous even function of period two
equal to \(1-|t|\) on \([-1,1]\). Then, uniformly in real \(t\),
\[
 \psi(t)=\frac12+
   \sum_{n\in\mathbb Z}
     \frac{2}{\pi^2(2n-1)^2}e^{i(2n-1)\pi t}.
\]
It is the characteristic function of the discrete law with
mass \(1/2\) at zero and mass
\(2/(\pi^2(2n-1)^2)\) at \((2n-1)\pi\), \(n\in\mathbb Z\).
This law differs from the Polya density law, although their
transforms agree throughout \([-1,1]\).
\end{hbtheorem}

\begin{hbtheorem}{thm:p15-criteria-equality-sets}{proved=annex:p15.criteria-equality-sets}{Every possible closed equality set}
For two characteristic functions, their equality set is
closed, contains zero and is symmetric under \(t\mapsto-t\).
Conversely, for every closed symmetric \(A\subseteq\mathbb R\)
containing zero, there is a characteristic function
\(\varphi_2\) with
\[
 \{t:e^{-|t|}=\varphi_2(t)\}=A.
\]
The annex gives an explicit \(C^2\) perturbation on the
components of \((0,\infty)\setminus A\), with derivative
bounds ensuring Polya's hypotheses.
This is the full construction in Exercise~3.3.21.
\end{hbtheorem}

\begin{hbtheorem}{thm:p15-criteria-cancellation}{proved=annex:p15.criteria-cancellation}{Both cancellation exercises}
There are mutually independent \(X,Y,Z\) with
\(Y\not\stackrel d=Z\) but \(X+Y\stackrel d=X+Z\):
take \(X,Y\) with the Polya density law and \(Z\) with the
preceding periodic discrete law.
On the other hand, if \(X,Y\) are independent and
\(X+Y\stackrel d=X\), then \(Y=0\) almost surely.
These are Exercises~3.3.22--3.3.23.
\end{hbtheorem}

\section{The moment problem and its countermodels}
\label{sec:p15-criteria-moment-problem}

\begin{hbdefinition}{def:p15-criteria-moment-problem}{Moments and the two determinacy domains}
A probability law with all absolute moments finite has moments
\(m_k=\int x^k\,\mu(dx)\), \(k\geq0\), with \(m_0=1\).
It is \emph{Hamburger determinate} if it is the only Borel
probability on \(\mathbb R\) with these moments.
A law supported on \([0,\infty)\) is \emph{Stieltjes determinate}
if it is unique among probability laws on that half-line.
The competing-law domain is part of each assertion.
For a real \(X\), write \(\nu_k=\mathbb E|X|^k\).
\index{moment problem}\index{moment determinacy}
\end{hbdefinition}

\begin{hbtheorem}{thm:p15-criteria-complex-gamma}{proved=annex:p15.criteria-complex-gamma}{The full complex gamma-integral premise}
For \(p>0\) and \(q=a+ib\) with \(a>0\), define
\[
 \arg q=\arctan(b/a),\qquad
 q^{-p}=\exp\{-p(\log|q|+i\arg q)\},\qquad
             -\pi/2<\arg q<\pi/2.
\]
Then the absolutely convergent integral satisfies
\[
 \int_0^\infty t^{p-1}e^{-qt}\,dt=\Gamma(p)q^{-p}.
\]
The proof uses a real parameter, integration by parts and
the preceding differentiation theorem, not a presumed
complex contour substitution.
\end{hbtheorem}

\begin{hbtheorem}{thm:p15-criteria-lognormal-family}{proved=annex:p15.criteria-lognormal-family}{Counterexample 1: all continuous and discrete source families}
For \(x>0\), let
\[
 f_0(x)=\frac{e^{-(\log x)^2/2}}{\sqrt{2\pi}\,x},
 \qquad
 f_a(x)=f_0(x)\{1+a\sin(2\pi\log x)\},\quad -1\leq a\leq1,
\]
and set the densities to zero for \(x\leq0\).
The laws indexed by distinct \(a\) are distinct, with moments
\(\int x^r f_a(x)dx=e^{r^2/2}\), \(r=0,1,\ldots\).
More generally
\[
 f_0(x)\left\{1+\sum_{k=1}^\infty a_k\sin(k\pi\log x)\right\}
 \quad(x>0),\qquad \sum_{k\geq1}|a_k|\leq1,
\]
gives probability densities with the same moments.
For every \(a>0\), the discrete law
\[
 \mathbb P(Y_a=ae^k)=\frac{a^{-k}e^{-k^2/2}}{c_a},
 \qquad c_a=\sum_{k\in\mathbb Z}a^{-k}e^{-k^2/2},
\]
is also normalized, has finite moments of every order, and
\(\mathbb EY_a^r=e^{r^2/2}\) for integers \(r\geq0\).
The normalization is finite and positive. Distinct parameters
in \([1,e)\) give disjoint discrete supports; parameters
differing by a factor \(e^j\), \(j\in\mathbb Z\), give
the same discrete law.
\end{hbtheorem}

\begin{hbtheorem}{thm:p15-criteria-real-counterexample}{proved=annex:p15.criteria-real-counterexample}{Counterexample 2: the whole real-line family}
Fix \(0<\lambda<1\), put \(\beta=\tan(\lambda\pi/2)\) and
\(c_\lambda=\lambda/(2\Gamma(1/\lambda))\).
For \(-1\leq a\leq1\),
\[
 f_{a,\lambda}(x)=c_\lambda e^{-|x|^\lambda}
       \{1+a\sin(\beta|x|^\lambda\operatorname{sgn}x)\}
                                      \quad(x\in\mathbb R)
\]
are distinct probability densities with all absolute moments
finite and the same moments for fixed \(\lambda\).
Their odd moments are zero and their even moments are
\(\Gamma((2k+1)/\lambda)/\Gamma(1/\lambda)\).
The full odd-integral cancellation uses the complex-gamma
identity with its specified branch.
\end{hbtheorem}

\begin{hbtheorem}{thm:p15-criteria-compact-moments}{proved=annex:p15.criteria-compact-moments}{Moment growth at the essential bound and compact determinacy}
For finite real \(X\), let
\[
 \Lambda=\sup\{x:\mathbb P(|X|<x)<1\}\in[0,\infty].
\]
Then \(\nu_k^{1/k}\to\Lambda\), with infinite values permitted.
Every compactly supported probability law is Hamburger
determinate. Its full proof first forces any competing law
with the same moments onto the same compact interval, then
uses uniform polynomial approximation.
This supplies the opening compact-support argument and
Exercise~3.3.24.
\end{hbtheorem}

\begin{hbtheorem}{thm:p15-criteria-analytic-determinacy}{proved=annex:p15.criteria-analytic-determinacy}{Durrett 3.3.25: the full analytic propagation proof}
For a prescribed finite moment sequence \(m_k\), if
\[
 \limsup_{k\to\infty}\frac{m_{2k}^{1/(2k)}}{2k}
                  =r<\infty,
\]
there is at most one probability law with those moments.
The expression is applied to nonnegative even moments;
if a positive-order even moment is zero the only possible
law is \(\delta_0\).
For any candidate law, all its characteristic derivatives
exist, and at every real center \(\theta\),
\[
 \varphi(\theta+t)
    =\sum_{j=0}^\infty \frac{t^j}{j!}\varphi^{(j)}(\theta)
                  \quad(|t|<1/(er))
\]
when \(r>0\). For \(r=0\) the series holds for every real
\(t\). The proof treats both cases and propagates equality
between centers explicitly.
\end{hbtheorem}

\begin{hblemma}{thm:p15-criteria-quasianalytic}{proved=annex:p15.criteria-quasianalytic}{An explicit quasi-analyticity lemma}
Let \(M_n>0\), \(n\geq0\), satisfy
\(M_n^2\leq M_{n-1}M_{n+1}\) for \(n\geq1\).
Suppose a complex \(C^\infty\) function \(f\) on \(\mathbb R\)
obeys \(|f^{(n)}(x)|\leq M_n\) for every \(x,n\), and
\[
 \sum_{n=1}^\infty (M_0/M_n)^{1/n}=\infty.
\]
If all derivatives of \(f\) vanish at zero, then \(f\equiv0\).
The proof supplies the weighted arithmetic--geometric
estimate and the derivative-envelope argument in full.
\end{hblemma}

\begin{hbtheorem}{thm:p15-criteria-carleman}{proved=annex:p15.criteria-carleman}{The complete Carleman criterion mentioned by the source}
A probability law with all finite moments is Hamburger
determinate if
\[
 \sum_{k=1}^\infty m_{2k}^{-1/(2k)}=\infty.
\]
A term with \(m_{2k}=0\) is interpreted as \(+\infty\),
and its zero-law case is treated separately.
The growth condition in Theorem~\ref{thm:p15-criteria-analytic-determinacy}
implies this condition.
The annex gives a complete independent proof of the
criterion stated without proof in the source's remark.
\index{Carleman's criterion}
\end{hbtheorem}

\begin{hbtheorem}{thm:p15-criteria-moment-method}{proved=annex:p15.criteria-moment-method}{Durrett 3.3.26: existence, moment transfer and weak convergence}
Suppose each \(\mu_n\) has all finite absolute moments,
\(\int x^k\,\mu_n(dx)\to m_k\in\mathbb R\) for every
positive integer \(k\), and
\[
 \limsup_{k\to\infty}m_{2k}^{1/(2k)}/(2k)<\infty.
\]
Then there exists a unique probability law \(\mu\) with
these moments and \(\mu_n\Rightarrow\mu\).
The proof establishes tightness and every transferred
moment before invoking uniqueness.
\end{hbtheorem}

\begin{hbtheorem}{thm:p15-criteria-generalized-gamma}{proved=annex:p15.criteria-generalized-gamma}{Exercise 3.3.25 with all parameter and growth conditions}
For \(\lambda>0\) and \(\alpha>-1\), the normalized density
of \(W=|X|\) proportional to \(x^\alpha e^{-x^\lambda}\)
on \((0,\infty)\) has
\[
 C=\frac{\lambda}{\Gamma((\alpha+1)/\lambda)},\qquad
 \nu_n=\frac{\Gamma((n+\alpha+1)/\lambda)}
                  {\Gamma((\alpha+1)/\lambda)}.
\]
Its exact root asymptotic is
\[
 \frac{\nu_n^{1/n}}n
       \sim(e\lambda)^{-1/\lambda}n^{1/\lambda-1}.
\]
Thus the source growth condition holds for \(\lambda\geq1\)
and fails for \(0<\lambda<1\).
In particular every normal law and every gamma law is
determined by its moments. Failure of this sufficient
condition alone does not assert indeterminacy for every
law in the displayed family.
\end{hbtheorem}

\section{Nonnegative laws and the last counterexample}
\label{sec:p15-criteria-stieltjes}

\begin{hbtheorem}{thm:p15-criteria-stieltjes}{proved=annex:p15.criteria-stieltjes}{The symmetric-square correspondence and the nonnegative criterion}
There is a bijection between probability laws on
\([0,\infty)\) and symmetric probability laws on \(\mathbb R\):
map \(X\geq0\) to \(Y=\varepsilon\sqrt X\), where
\(\varepsilon\) is an independent fair sign, and recover
the first law by \(Y^2\).
If \(X\) has all finite moments \(\nu_k\), then
\(\mathbb EY^{2k}=\nu_k\) and \(\mathbb EY^{2k+1}=0\).
Consequently
\[
 \limsup_{k\to\infty}\nu_k^{1/(2k)}/(2k)<\infty
\]
implies Stieltjes determinacy.
For the preceding generalized-gamma family on the
nonnegative half-line, this sufficient condition holds
when \(\lambda\geq1/2\).
\end{hbtheorem}

\begin{hbtheorem}{thm:p15-criteria-positive-counterexample}{proved=annex:p15.criteria-positive-counterexample}{Counterexample 3: the complete nonnegative family}
Fix \(0<\lambda<1/2\), set
\(\beta=\tan(\lambda\pi)\) and
\(c_\lambda=\lambda/\Gamma(1/\lambda)\).
For \(-1\leq a\leq1\), the functions
\[
 f_a(x)=c_\lambda e^{-x^\lambda}
                     \{1+a\sin(\beta x^\lambda)\}
                    \quad(x\geq0),
\]
and zero for \(x<0\), are distinct probability densities.
All have the same finite moments
\(\Gamma((n+1)/\lambda)/\Gamma(1/\lambda)\), \(n\geq0\).
This proves the source's final counterexample, including
its normalization and every integer moment.
\end{hbtheorem}

\begin{proofboundary}{def:p15-criteria-completion-boundary}{Source completion and the next presentation}
Together with Chapter~\ref{sec:p15-characteristic-functions-and-inversion},
this chapter supplies Durrett Section~3.3's numbered
theorems and examples, its three moment-problem counterexamples
and Exercises~3.3.1--3.3.25.
The Carleman remark has received its own full proof.
The i.i.d. and finite-row central limit theorems, both array routes,
all Section~3.4 exercises, Erd\H{o}s--Kac and Berry--Esseen are supplied
in \cref{sec:p15-central-limit-theorems}.
The complete normal-domain characterization and the original
truncated-mean and cutoff route are supplied in
\cref{thm:p15-nda-characterization,thm:p15-nda-gk-cutoff-route}.
The later full stable-law presentations retain their distinct
constructions.
No CLT was assumed in the present chapter.
\end{proofboundary}
